\section{Rigidity of the first-letter identities}\label{sec:rigidity}

The identities furnished by the walk concern words whose unsigned indices
are pairwise distinct.  We prove that their two left multipliers must
agree.  This rules out the distinct martingale endpoints constructed
in Section~\ref{sec:liouville}.

\begin{theorem}\label{thm:rigidity}
Let $(N,\tau)$ be a finite von Neumann algebra with a faithful normal
tracial state.  Let $R:N\to N$ be bounded and complex linear, and suppose
that
\begin{equation}\label{eq:rigidity-continuity}
 \norm{R(x_n)}_2\longrightarrow0
 \quad\text{whenever}\quad
 \sup_n\norm{x_n}<\infty,
 \qquad \norm{x_n}_2\longrightarrow0.
\end{equation}
Let $(g_j)_{j\geq1}$ be a free Haar family in $N$, and let $a,b\in N$.
Suppose that for every nonempty signed word
$w=g_{i_1}^{\varepsilon_1}\cdots g_{i_q}^{\varepsilon_q}$ with
$\varepsilon_j\in\{1,-1\}$ and pairwise distinct indices $i_1,\ldots,i_q$,
\begin{equation}\label{eq:rigidity-first-letter}
 R(w)=
 \begin{cases}
 aw,&\varepsilon_1=1,\\
 bw,&\varepsilon_1=-1.
 \end{cases}
\end{equation}
Then $a=b$.
\end{theorem}

Here the free Haar condition means that the trace of every nonidentity
reduced group word in the $g_j$ is zero.  In particular, distinct reduced
words are orthonormal for the tracial $2$-norm.  The coefficients $a,b$
are arbitrary elements of the containing algebra.

The target is a unitary whose positive powers carry the multiplier $a$
and whose negative powers carry the multiplier $b$.  A spectral argument
will show that these identities force $a=b$ for a bounded map in a finite
tracial algebra.  Higher powers of any one $g_j$ repeat its unsigned index, so
the hypothesis does not give this target directly.  Instead, we average
many free generators, control the repeated indices in suitable polynomial
expressions, and pass to the polar part of the resulting element in a
further tracial ultrapower.

\subsection{Norm control for free words}

We use the following form of Haagerup's length estimate
\cite[Lemma~1.4]{Haagerup1979}.  We give the proof, including the transfer
from the regular representation to the containing tracial algebra.

\begin{lemma}\label{lem:rigidity-length}
Let $\mathbb F_r$ be the free group on $r$ generators and let $\lambda$
be its left regular representation.  For a finitely supported family
$(c_w)_{|w|=q}$ on the reduced words of length $q\geq0$,
\begin{equation}\label{eq:rigidity-length}
 \left\|\sum_{|w|=q}c_w\lambda(w)\right\|
 \leq(q+1)\left(\sum_{|w|=q}|c_w|^2\right)^{1/2}.
\end{equation}
If $g_1,\ldots,g_r$ have free Haar joint moments in a faithfully tracial
von Neumann algebra, every scalar-coefficient word polynomial in them
has the same norm as the corresponding polynomial in the regular
generators.  Consequently, \eqref{eq:rigidity-length} also holds with
$\lambda(w)$ replaced by $w(g)$, and its right-hand side is
$(q+1)\norm{\sum_{|w|=q}c_w w(g)}_2$.
\end{lemma}

\begin{proof}
Write $L=\sum_{|w|=q}c_w\lambda(w)$.  On the standard basis of
$\ell^2(\mathbb F_r)$, separate each product $wx$ according to its
exact number $t\in\{0,\ldots,q\}$ of cancellations at the joining
boundary.  This gives $L=\sum_{t=0}^q L_t$ on finitely supported vectors.
Fix $t$.  Such a product has the form
\[
 w=uv,\qquad x=v^{-1}h,\qquad wx=uh,
 \qquad |u|=q-t,\quad |v|=t,
\]
where all three displayed concatenations $uv$, $v^{-1}h$, and $uh$
are reduced.  Conversely, these conditions describe exactly the
products with $t$ cancellations, allowing any of $u,v,h$ to be empty
when its required length permits this.  In particular, $h$ may be empty
when the input is exhausted.

For this fixed $t$, an input basis word of length at least $t$ has a
unique prefix of length $t$ and hence a unique suffix $h$.  An output
basis word of length at least $q-t$ likewise determines $h$ uniquely by
removing its prefix of length $q-t$.  Thus the input spaces indexed by
different $h$ are orthogonal, as are the corresponding output spaces.
Relative to these two decompositions, $L_t$ is a direct sum of matrices
$B_h$, with columns indexed by $v$ and rows by $u$.  An entry is
$c_{uv}$ when the three reduction conditions hold and is zero otherwise.
The fixed cutting length makes $(u,v)\mapsto uv$ injective within each
matrix.  Therefore
\[
 \norm{B_h}\leq\norm{B_h}_{\mathrm{HS}}
 \leq\left(\sum_{|w|=q}|c_w|^2\right)^{1/2}.
\]
Taking the direct-sum norm bounds $L_t$ by the same quantity.  Summing
over the $q+1$ possible values of $t$ proves
\eqref{eq:rigidity-length}.

For the norm transfer, recall that if $A\geq0$ belongs to a von Neumann
algebra with faithful tracial state $\rho$, then
\begin{equation}\label{eq:rigidity-norm-moments}
 \norm{A}=\lim_{n\to\infty}\rho(A^n)^{1/n}.
\end{equation}
The upper bound is immediate.  If $0<c<\norm{A}$, the spectral
projection $e=1_{(c,\infty)}(A)$ is nonzero and has positive trace, so
$\rho(A^n)\geq c^n\rho(e)$.  Taking roots and then letting
$c\uparrow\norm{A}$ proves the formula.  For a word polynomial $p$,
the free Haar moments give, for every $n\geq1$,
\[
 \tau\bigl((p(g)^*p(g))^n\bigr)
 =\tau_{\mathbb F_r}\bigl((p(\lambda)^*p(\lambda))^n\bigr).
\]
Both traces are faithful; here $\tau_{\mathbb F_r}$ is the canonical
trace on the free group von Neumann algebra.  Applying
\eqref{eq:rigidity-norm-moments} gives $\norm{p(g)}=\norm{p(\lambda)}$.
The asserted $2$-norm formula follows from orthonormality of reduced
words.
\end{proof}

Under the hypotheses of Theorem~\ref{thm:rigidity}, set
\begin{equation}\label{eq:rigidity-normalized-sums}
 s_r=r^{-1/2}\sum_{j=1}^r g_j\in N.
\end{equation}
Lemma~\ref{lem:rigidity-length} gives $\norm{s_r}_2=1$ and
$\norm{s_r}\leq2$ for every $r$.  Fix a free ultrafilter $\omega_2$ on
$\N$ and put
\begin{equation}\label{eq:rigidity-second-ultrapower}
 \widehat N=N^{\omega_2},\qquad
 \widehat\tau=\tau^{\omega_2},\qquad
 s=(s_r)_{\omega_2}\in\widehat N,\qquad S=R^{\omega_2}.
\end{equation}
Lemma~\ref{lem:ultrapower} ensures that $S$ is bounded and retains
$2$-norm continuity on operator norm balls.  We identify $a,b$ with
their constant images in $\widehat N$.

In the application from Section~\ref{sec:liouville},
$N=M^{\omega_1}$ and $\widehat N=(M^{\omega_1})^{\omega_2}$.
The first index $m$ selects increasingly accurate unitary tuples; the
second index $r$ increases the number of exact free generators being
averaged in $N$, where the bound $\norm{s_r}\leq2$ is already available.
The hypotheses and notation of Theorem~\ref{thm:rigidity} remain in
force through the next two subsections.

Figure~\ref{fig:proof-map} locates the two ultrapowers in the Liouville
argument and records the remaining passage from these bounded sums to a
polar unitary and then to $a=b$.
\begin{figure}[htbp]
\centering
\small
\begin{tabular}{c}
\fbox{\parbox{0.86\linewidth}{\centering
The fixed walk has asymptotically free moments.\\
A hypothetical nonconstant harmonic-map martingale has\\
distinct support points $a,b\in M$.}}\\[3pt]
$\downarrow$\quad finite selection and the first ultrapower\\[3pt]
\fbox{\parbox{0.86\linewidth}{\centering
In $N=M^{\omega_1}$: exact free Haar unitaries $g_1,g_2,\ldots$\\
and exact first-letter identities for $R=T^{\omega_1}$.}}\\[3pt]
$\downarrow$\quad free-group length estimate in $N$\\[3pt]
\fbox{\parbox{0.86\linewidth}{\centering
$s_r=r^{-1/2}\sum_{j=1}^r g_j$ satisfies $\norm{s_r}\leq2$.\\
In each fixed product of blocks $s_r(s_r^*s_r)^d$, $d\geq0$,\\
terms with repeated surviving indices have total norm tending to zero.}}\\[3pt]
$\downarrow$\quad the second ultrapower and polar approximation\\[3pt]
\fbox{\parbox{0.86\linewidth}{\centering
In $\widehat N=N^{\omega_2}$: $s=(s_r)_{\omega_2}$ has a unitary
polar part $V$,\\
and $S=R^{\omega_2}$ satisfies $S(V^k)=aV^k$, $S(V^{-k})=bV^{-k}$.}}\\[3pt]
$\downarrow$\quad bounded sine sums and the finite faithful trace\\[3pt]
\fbox{\parbox{0.86\linewidth}{\centering
$a=b$, a contradiction. Hence $T(x)=T(1)x$.\\
Last-input averaging then produces the tracial cocycle primitive.}}
\end{tabular}
\caption{The Liouville contradiction and its use for cocycles.
The first ultrapower uses scalar moment convergence; the norm bound
is proved for the exact free family before the second ultrapower.
The constants $a,b$ are carried through the trace-preserving constant
embeddings, and the power identities hold for every $k\geq1$.}
\label{fig:proof-map}
\end{figure}


\subsection{From distinct indices to ordered blocks}

Polynomial expressions in $s_r$ may contain repeated indices, so their
multiplier identities do not follow directly from the hypothesis.
For polar approximation we will need powers of $sp(s^*s)$, where $p$
is a scalar polynomial.  Expanding these powers without commuting any
factors gives products of the blocks $B_d(z)=z(z^*z)^d$, $d\geq0$.
We establish their identities by showing that repeated surviving indices
contribute an error tending to zero in operator norm.

\begin{lemma}\label{lem:rigidity-blocks}
For every $k\geq1$ and $d_1,\ldots,d_k\geq0$, let
\[
 P(z)=B_{d_1}(z)B_{d_2}(z)\cdots B_{d_k}(z).
\]
Then
\begin{equation}\label{eq:rigidity-block-identities}
 S(P(s))=aP(s),\qquad S(P(s)^*)=bP(s)^*.
\end{equation}
\end{lemma}

\begin{proof}
Fix this product, and let $L=\sum_{j=1}^k(2d_j+1)$ be its degree.
Every term in the expansion of one block $B_d(s_r)$ has the odd
alternating sign pattern
\[
 +,-,+,-,\ldots,+.
\]
Deleting adjacent inverse letters deletes two consecutive opposite
signs.  The remaining pattern is still alternating, begins and ends
positively, and is nonempty: the excess of positive over negative
signs remains one.  Hence the free reduction of each block is a
nonempty word with positive endpoint signs.  After reducing the blocks
internally, their boundaries join two positive letters and admit no
further cancellation.  Every reduced word in $P(s_r)$ is therefore
nonempty and begins and ends positively.  Taking adjoints reverses
the block order and all signs, so the corresponding assertion for
$P(s_r)^*$ has negative endpoint signs.

We next bound the coefficient of every possible reduced word, including
words much shorter than $L$.  Write
\[
 P(s_r)=\sum_{q=1}^L\sum_{|w|=q}c_w(r)w(g).
\]
For a fixed reduced word $w$ of length $q$, record a free reduction of
an original term by pairing its original positions whenever their
letters cancel.  For definiteness, the reduction may be performed by
reading from left to right and canceling against the last surviving
letter whenever possible.  The recorded pairs form a noncrossing
partial matching of the $L$ positions: there are $(L-q)/2$ pairs, and
the unmatched positions spell $w$ in order.  All nested and newly
adjacent cancellations are included in this record.

There are at most $3^L$ noncrossing partial matchings of $L$ positions.
Indeed, mark each position as an opening endpoint, a closing endpoint,
or unmatched; noncrossing forces a closing endpoint to pair with the
last unpaired opening endpoint, so this encoding is injective.  We
may count every such partial matching, including ones not realized by
a reduction.  For a fixed pattern and fixed $w$, the labels at the
unmatched positions are prescribed.  Each canceled pair has equal
unsigned labels and opposite signs and hence has at most $r$ label
choices.  Coincidences between different pairs or with unmatched
positions can only reduce the number of choices.  Each original term
has coefficient $r^{-L/2}$.  Consequently, for every $q$ and $w$,
\begin{equation}\label{eq:rigidity-coefficient-bound}
 |c_w(r)|\leq C_L r^{(L-q)/2}r^{-L/2}
 =C_Lr^{-q/2},\qquad C_L=3^L.
\end{equation}
An impossible parity or reduction pattern contributes zero.  The same
bound applies to the adjoint polynomial.

Let $E_{r,q}$ be the sum of those length-$q$ terms whose \emph{surviving
reduced word} has a repeated unsigned index.  Original repetitions
which disappear in the reduction remain included in the coefficients
of the retained words.  For $q=1$ this sum is zero.  For $q\geq2$ the
number of possible words in the sum is at most
\[
 2^q\binom q2 r^{q-1}:
\]
choose a pair of equal-index positions, choose their common label and
the other $q-2$ labels, and choose all signs.  Counting nonreduced
strings here only enlarges the bound.  By orthonormality and
\eqref{eq:rigidity-coefficient-bound},
\begin{equation}\label{eq:rigidity-discard-ltwo}
 \norm{E_{r,q}}_2^2
 \leq C_L^2 2^q\binom q2 r^{-1}\qquad(2\leq q\leq L).
\end{equation}
Lemma~\ref{lem:rigidity-length} now yields the aggregate operator norm
estimate
\begin{equation}\label{eq:rigidity-discard-norm}
 E_r=\sum_{q=1}^L E_{r,q},\qquad
 \norm{E_r}\leq\sum_{q=2}^L(q+1)\norm{E_{r,q}}_2
 =O_L(r^{-1/2})\longrightarrow0.
\end{equation}
There is no scalar term, by the block reduction argument.

Put $Q_r=P(s_r)-E_r$.  Each word in this finite sum has pairwise
distinct unsigned indices and begins positively.  Thus
\eqref{eq:rigidity-first-letter} and linearity give $R(Q_r)=aQ_r$.
Moreover, $\norm{P(s_r)}\leq2^L$, so \eqref{eq:rigidity-discard-norm}
also proves that the sequence $(Q_r)$ is bounded.  In $N$ we have
\[
 \norm{R(P(s_r))-aP(s_r)}
 \leq(\norm{R}+\norm{a})\norm{E_r}\longrightarrow0.
\]
Passing to $\widehat N$ proves the first identity in
\eqref{eq:rigidity-block-identities}.  For the adjoint, the retained
sum is $Q_r^*$: its reduced words have distinct unsigned indices and
begin negatively.  Therefore $R(Q_r^*)=bQ_r^*$, and the identical
estimate with $b$ proves the second identity.
\end{proof}

\subsection{A polar unitary and all its powers}

The spectral calculation below is related to the polar decomposition
of circular elements in free probability, established by Voiculescu
and also proved by Banica~\cite{Voiculescu1990Circular,Banica1996Polar}.
We derive the required Catalan moments and spectral distribution
directly. Only the absence of an atom at zero is needed to obtain
the unitary polar part.

\begin{lemma}\label{lem:rigidity-polar}
The polar part $V$ of $s$ is a unitary in $\widehat N$, and
\begin{equation}\label{eq:rigidity-power-identities}
 S(V^k)=aV^k,\qquad S(V^{-k})=bV^{-k}
 \qquad(k\geq1).
\end{equation}
\end{lemma}

\begin{proof}
We first show that $s$ has zero kernel, which makes its polar part
unitary in the finite algebra.  We then transfer the block identities
to powers of that unitary.
Set $A=s^*s$, so $0\leq A\leq4$.  To determine its kernel, we compute
its moments.  For
$d\geq1$, the quantity $\tau((s_r^*s_r)^d)$ is $r^{-d}$ times the
number of assignments of labels in $\{1,\ldots,r\}$ to the alternating
$2d$ positions
\[
 -,+,-,+,\ldots,-,+
\]
whose word reduces to the identity.  An identity reduction pairs all
positions noncrossingly, so it uses at most $d$ unsigned labels.
Assignments using at most $d-1$ labels number $O_d(r^{d-1})$: for each
number of labels, choose an ordered list of those labels and then
assign a member of that list to each of the $2d$ positions.

If exactly $d$ labels occur in an identity word, every label occurs
twice.  Its two occurrences must be paired in any cancellation
matching, so the matching is uniquely determined by the labeling.
Conversely, in a noncrossing perfect matching, every pair has an even
number of positions inside it, since those positions must be matched
among themselves.  The endpoints therefore have opposite parities,
and hence opposite signs in our alternating string.  Assigning
distinct labels to its $d$ pairs gives an identity word: erase
innermost pairs successively.  It follows that the number of leading
assignments is exactly
\[
 C_d(r)_d,\qquad (r)_d=r(r-1)\cdots(r-d+1),
\]
where $C_d$ is the number of noncrossing perfect matchings on $2d$
positions.  Choosing the partner of the first position gives
$C_0=1$ and $C_d=\sum_{j=0}^{d-1}C_jC_{d-1-j}$.  The generating
function thus satisfies $C(z)=1+zC(z)^2$; its solution with constant
term one gives
\[
 C_d=\frac1{d+1}\binom{2d}{d}.
\]
Consequently,
\begin{equation}\label{eq:rigidity-catalan-moments}
 \widehat\tau(A^d)
 =\lim_{r\to\omega_2}\tau((s_r^*s_r)^d)
 =\lim_{r\to\infty}r^{-d}
       \bigl(C_d(r)_d+O_d(r^{d-1})\bigr)
 =C_d.
\end{equation}
The zeroth moment is also $C_0=1$.

Consider the measure
\begin{equation}\label{eq:rigidity-mp-law}
 d\nu(x)=\frac1{2\pi}\sqrt{\frac{4-x}{x}}\,
                1_{(0,4)}(x)\,dx.
\end{equation}
For every integer $d\geq0$, substitution $x=4t$ and the beta integral
give
\[
 \begin{split}
 \int_0^4 x^d\,d\nu(x)
 &=\frac{4^{d+1}}{2\pi}
     \int_0^1 t^{d-1/2}(1-t)^{1/2}\,dt\\
 &=\frac{4^{d+1}}{2\pi}
     \frac{\Gamma(d+\tfrac12)\Gamma(\tfrac32)}{\Gamma(d+2)}
 =\frac{(2d)!}{d!(d+1)!}=C_d.
 \end{split}
\]
In particular, $\nu$ is a probability measure.  Both $\nu$ and the
spectral distribution of $A$ are supported in $[0,4]$.  Equality of
all their moments implies equality on polynomials and then on all
continuous functions by uniform polynomial approximation.  They are
therefore the same measure.  Since \eqref{eq:rigidity-mp-law} has no
atom at zero,
\[
 \widehat\tau(1_{\{0\}}(A))=0,
 \qquad 1_{\{0\}}(A)=0
\]
by faithfulness.  In the polar decomposition $s=VA^{1/2}$ this gives
$V^*V=1$.  Traciality and faithfulness then give
\[
 \widehat\tau(1-VV^*)=1-\widehat\tau(V^*V)=0,
 \qquad VV^*=1.
\]
Thus $V$ is a unitary.

It remains to transfer the identities.  For a fixed regularization
parameter we will use polynomial approximation in operator norm;
we will then remove the regularization using bounded-ball $2$-continuity.
For $\varepsilon>0$ put
\[
 x_\varepsilon=s(A+\varepsilon1)^{-1/2}
 =V\left(\frac{A}{A+\varepsilon1}\right)^{1/2}.
\]
Fix $k\geq1$.  A polynomial $p(t)=\sum_{d=0}^D c_dt^d$ has the
literal ordered expansion
\begin{equation}\label{eq:rigidity-ordered-expansion}
 [sp(A)]^k
 =\sum_{d_1,\ldots,d_k=0}^D c_{d_1}\cdots c_{d_k}
       [sA^{d_1}]\cdots[sA^{d_k}].
\end{equation}
Every summand is a product covered by Lemma~\ref{lem:rigidity-blocks}.
That lemma and complex linearity yield
\[
 S([sp(A)]^k)=a[sp(A)]^k,
 \qquad S(([sp(A)]^k)^*)=b([sp(A)]^k)^*.
\]
For the second identity one takes the adjoint of the expansion
\eqref{eq:rigidity-ordered-expansion}, conjugating its scalar
coefficients and reversing the block order, and applies the separate
adjoint identity in Lemma~\ref{lem:rigidity-blocks}.
Choose polynomials converging uniformly on $[0,4]$ to
$(t+\varepsilon)^{-1/2}$.  Functional calculus and boundedness of $S$
pass the last two identities to the norm limit, giving
\begin{equation}\label{eq:rigidity-regularized-powers}
 S(x_\varepsilon^k)=ax_\varepsilon^k,
 \qquad S((x_\varepsilon^k)^*)=b(x_\varepsilon^k)^*.
\end{equation}

The $x_\varepsilon$ are contractions, and their distance to $V$
satisfies
\[
 \norm{x_\varepsilon-V}_2^2
 =\int_0^4
   \left(\sqrt{\frac{x}{x+\varepsilon}}-1\right)^2\,d\nu(x)
 \longrightarrow0
\]
by dominated convergence and $\nu(\{0\})=0$.  Multiplication satisfies
$\norm{xyz}_2\leq\norm{x}\norm{y}_2\norm{z}$ in a tracial algebra.
The noncommutative telescoping identity
\[
 x_\varepsilon^k-V^k
 =\sum_{j=0}^{k-1}x_\varepsilon^j
                  (x_\varepsilon-V)V^{k-1-j}
\]
therefore shows that
$\norm{x_\varepsilon^k-V^k}_2\leq k\norm{x_\varepsilon-V}_2\to0$;
the same holds for their adjoints.  All these powers are contractions.
Using the bounded-ball $2$-norm continuity of $S$ in
\eqref{eq:rigidity-regularized-powers}, along for example
$\varepsilon=1/n$, proves \eqref{eq:rigidity-power-identities}.
\end{proof}

\subsection{The spectral conclusion}

The final step uses boundedness, the identities for powers, and the
finite faithful trace. No continuity assumption on the map is needed
in Lemma~\ref{lem:rigidity-unitary}.

\begin{lemma}\label{lem:rigidity-unitary}
Let $(\mathcal A,\rho)$ be a finite von Neumann algebra with faithful
normal tracial state, let $V\in\mathcal A$ be unitary, and let
$S:\mathcal A\to\mathcal A$ be bounded and complex linear.  If
$a,b\in\mathcal A$ satisfy
\[
 S(V^k)=aV^k,\qquad S(V^{-k})=bV^{-k}\qquad(k\geq1),
\]
then $a=b$.
\end{lemma}

\begin{proof}
We use analytic polynomials whose imaginary parts remain uniformly
bounded on the unit circle while their real parts grow near $1$.
Rotating these polynomials
will bound $a-b$ on each short spectral arc; the trace will then combine
the compression estimates without a loss depending on the number of arcs.
Write $\mathbb T=\{z\in\C:|z|=1\}$ and
\[
 H_m(z)=\sum_{k=1}^m\frac{z^k}{k},\qquad
 h_m=\sum_{k=1}^m\frac1k.
\]
We first verify the uniform sine-sum bound
\begin{equation}\label{eq:rigidity-sine-bound}
 \sup_{m\geq1}\sup_{\theta\in\R}
 \left|\sum_{k=1}^m\frac{\sin(k\theta)}k\right|
 \leq B,\qquad B=1+\pi.
\end{equation}
By periodicity and oddness it suffices to use the angular distance
$t\in[0,\pi]$ from $2\pi\Z$.  At $t=0$ the sum is zero.  For $t>0$
put $K=\lfloor1/t\rfloor$.  The part with $k\leq\min(K,m)$ has
absolute sum at most $Kt\leq1$, since $|\sin(kt)|\leq kt$.
Geometric summation gives, for every interval of integers $p\leq q$,
\[
 \left|\sum_{k=p}^q\sin(kt)\right|
 \leq\left|\sum_{k=p}^q e^{ikt}\right|
 \leq\frac1{\sin(t/2)}\leq\frac\pi t.
\]
If $m\geq K+1$, put $p=K+1$ and
$A_j=\sum_{k=p}^j\sin(kt)$.  Abel summation gives
\[
 \sum_{k=p}^m\frac{\sin(kt)}k
 =\frac{A_m}{m}
   +\sum_{j=p}^{m-1}A_j\left(\frac1j-\frac1{j+1}\right).
\]
Its absolute value is at most $\pi/(tp)\leq\pi$, including when
$K=0$.  This proves \eqref{eq:rigidity-sine-bound}.

For $\lambda\in\mathbb T$, write
\[
 H_m(\lambda V)=X_{m,\lambda}+iY_{m,\lambda},
 \qquad X_{m,\lambda}=X_{m,\lambda}^*,\quad
 Y_{m,\lambda}=Y_{m,\lambda}^*.
\]
Continuous functional calculus and \eqref{eq:rigidity-sine-bound}
give $\norm{Y_{m,\lambda}}\leq B$.  The identities for positive and
negative powers imply
\[
 \begin{split}
 S\bigl(H_m(\lambda V)-H_m(\lambda V)^*\bigr)
 &=aH_m(\lambda V)-bH_m(\lambda V)^*\\
 &=(a-b)X_{m,\lambda}+i(a+b)Y_{m,\lambda}.
 \end{split}
\]
Since the argument of $S$ has norm at most $2B$, we obtain
\begin{equation}\label{eq:rigidity-real-bound}
 \norm{(a-b)X_{m,\lambda}}\leq C,
 \qquad C=B(2\norm{S}+\norm{a}+\norm{b}),
\end{equation}
uniformly in $m$ and $\lambda$.

Fix $m$ and any half-open arc $I\subset\mathbb T$ of angular length
at most $1/m$, and let $e=1_I(V)$.  Choose $\lambda$ to rotate the
midpoint of $I$ to $1$.  On the rotated arc the argument $\theta$
satisfies $|\theta|\leq1/m$, so
\[
 \operatorname{Re}H_m(e^{i\theta})
 =\sum_{k=1}^m\frac{\cos(k\theta)}k
 \geq\cos(1)h_m.
\]
Thus, for $X=X_{m,\lambda}$, Borel functional calculus gives an
element $F\in W^*(V)$ equal to $1/X$ on $I$ and to zero off $I$,
with
\[
 XF=e,\qquad
 \norm{F}\leq\frac1{\cos(1)h_m}.
\]
Writing $c=a-b$ and multiplying \eqref{eq:rigidity-real-bound} on
the right by $F$ gives
\begin{equation}\label{eq:rigidity-arc-bound}
 \norm{ce}=\norm{cXF}
 \leq\frac{C}{\cos(1)h_m}=:\delta_m.
\end{equation}
In the corner $e\mathcal A e$, this says
$ec^*ce\leq\delta_m^2 e$.

Partition $\mathbb T$ into finitely many disjoint half-open arcs of
angular length at most $1/m$, and denote their spectral projections
by $e_1,\ldots,e_{J_m}$.  They are orthogonal and sum to $1$.
Traciality and \eqref{eq:rigidity-arc-bound} imply
\[
 \begin{split}
 \norm{c}_2^2
 &=\sum_{j=1}^{J_m}\rho(c^*ce_j)
  =\sum_{j=1}^{J_m}\rho(e_jc^*ce_j)\\
 &\leq\delta_m^2\sum_{j=1}^{J_m}\rho(e_j)
  =\delta_m^2\longrightarrow0,
 \end{split}
\]
because $h_m\geq\log(m+1)\to\infty$.  Each arc estimate is weighted
by the trace of its own projection; the number of arcs introduces no
factor.  Faithfulness gives $c=0$, as required.  The right multiplication
by $F$ and the tracial compression identity apply to arbitrary $a,b$;
the argument imposes no commutation condition and uses no particular
spectral distribution of $V$.
\end{proof}

\begin{proof}[Proof of Theorem~\ref{thm:rigidity}]
Construct $s\in\widehat N$ and $S=R^{\omega_2}$ as in
\eqref{eq:rigidity-second-ultrapower}.
Lemma~\ref{lem:rigidity-blocks} propagates the first-letter identities
to every ordered block product.  Lemma~\ref{lem:rigidity-polar} then
provides a unitary $V$ satisfying the positive and negative power
identities.  Lemma~\ref{lem:rigidity-unitary} gives equality of the
constant images of $a$ and $b$ in $\widehat N$.  The constant embedding
is trace preserving and faithful, so $a=b$ in $N$.
\end{proof}

\begin{remark}\label{rem:finite-trace-essential}
The finite hypothesis in Lemma~\ref{lem:rigidity-unitary} is essential.
On $H=\ell^2(\Z)$ let $Ve_n=e_{n+1}$, let $p$ project onto $\C e_0$,
and let $q$ project onto the closed span of the $e_n$ with $n<0$.
The contraction $S(X)=pXq$ on $B(H)$ satisfies
$S(V^k)=pV^k$ and $S(V^{-k})=0$ for every $k\ge1$:
the only nonzero input coordinate of $pV^k$ is $e_{-k}$, whereas
that of $pV^{-k}$ is $e_k$.
Thus the power identities hold with $a=p\ne0=b$.
\end{remark}
